\documentclass[12pt, a4paper]{article}

\usepackage[utf8]{inputenc}
\usepackage{amsfonts, amsmath, amssymb, amsthm}
\usepackage{geometry}
\usepackage[hidelinks]{hyperref}
\usepackage[none]{hyphenat}
\usepackage{setspace}
\usepackage{float}
\usepackage{fancyhdr}
\usepackage{enumitem}
\usepackage{graphicx}
\usepackage{cleveref}
\usepackage{tikz-cd}
\usepackage{mathrsfs}

\geometry{a4paper, left=0.5in, top=0.5in, right=0.5in, bottom=0.7in}
\onehalfspacing

\newtheorem{theorem}{Theorem}[section]
\newtheorem{lemma}{Lemma}[section]
\newtheorem{corollary}{Corollary}[theorem]
\theoremstyle{definition}
\newtheorem{definition}{Definition}[section]
\newtheorem{example}{Example}[section]
\theoremstyle{remark}
\newtheorem*{remark}{\textbf{Remark}}

\newcommand{\e}{\varepsilon}
\newcommand{\st}{\text{ such that }}
\newcommand{\B}{\mathscr{B}}
\DeclareMathOperator{\Span}{span}
\DeclareMathOperator{\nullity}{nullity}
\DeclareMathOperator{\rank}{rank}

\hyphenpenalty=10000
\exhyphenpenalty=10000

\title{\textbf{\Large Fundamentals of Matrix Algebra \\ \large Elementary Operations and the Mechanics of Rank}}
\author{\textbf{Tushar Kumar Aich}}
\date{\today}

\begin{document}
\maketitle

\begin{abstract}
    This paper explores the elementary matrix operations and then we define rank of a matrix and how these elementary operations are rank-invariant. We establish the canonical block-diagonal rank form via mathematical induction and understand that rank of a matrix is both linearly independent columns as well as rows of a matrix and that every invertible matrix can be created by products of elementary matrices. We conclude by proving the mechanics underlying the standard augmented matrix algorithm used for computing matrix inverses.
\end{abstract}

\section{Introduction}
A central goal in linear algebra is simplifying the matrix representation of a linear transformation without altering its core geometric invariants. The foundational link between these computational manipulations and abstract geometry is matrix rank, which measures the structural information contained within a system. By utilizing elementary row and column operations—which algebraically mirror coordinate basis changes—we can systematically reduce a matrix to an exceptionally clean skeleton. 

\section{Elementary Matrix Operations}
\begin{definition}
    Let $A$ be any $m \times n$ matrix. Any one of the following three operations on the rows/columns of $A$ is called an elementary row/column operation:
    \begin{enumerate}
        \item Interchanging any two rows/columns of $A$.
        \item Multiplying any row//column by a nonzero scalar.
        \item Adding any scalar multiple of a row/column to another row/column.
    \end{enumerate}
\end{definition}

Any of these operatoons is called an elementary operation. Elementary operations are of type 1, type 2 or type 3 depending on whether they are obtained by 1, 2 or 3.

\begin{definition}
    A $n \times n$ elementary matrix is a matrix obtained by performing an elementary operation on $I_n$. The elementary matrix is said to be of type 1, 2 or 3 according to whether the elementary row performed on $I_n$ is a type 1, 2 or 3 operation, respectively.
\end{definition}

So we can say that for any $m \times n$ matrix $A$ there always exists an $m \times m$ elementary matrix $E$ which is determined from $I_m$ which is used to perform a row operation on the matrix $A$ by multiplying the elementary matrix $E$ with $A$ to obtain the transformed matrix say $B$, \textit{i.e.}, $B = EA$.

Similarly there always exists an $n \times n$ elementary matrix $E$ which is determined from $I_n$ which is used to perform a column operation on the matrix $A$ by multiplying the elementary matrix $E$ with $A$ to obtain the transformed matrix say $B$, \textit{i.e.}, $B = AE$.

\begin{theorem}
    Elementary matrices are invertible, and the inverse of an elementary matrix is an elementary matrix of the same type.
\end{theorem}

\begin{proof}
    Let $E$ be an elementary matrix of order $n \times n$. Then $E$ can be obtained by performing an elementary row operation on $I_n$. We can transform $E$ back to $I_n$ by just reversing the process. Thus $I_n$ can be obtained from $E$ by an elementary row operation of same type. Thus there exists another elementary matrix $\overline{E}$ such that $E\overline{E} = I_n$. Hence $E$ is invertible and $E^{-1} = \overline{E}$.

    Same path is followed for proving through elementary column operations
\end{proof}

\section{Rank of a Matrix}
\begin{definition}
    If $A \in M_{m \times n}(F)$, we define the rank of $A$, to be the rank of the linear transformation $L_A:F^n \to F^m$.
\end{definition}

Here, we see that for a matrix of order $n \times n$, we define $\rank(A)$ to be $\rank(L_A)$ where $L_A:F^n \to F^m$. We know $A$ is invertible if and only if $L_A$ is invertible. Also we know that for two vector spaces of equal dimensions a transformation $T:V \to W$ is invertible if and only if $\rank(T) = \dim(V)$. Thus $L_A$ is invertible if and only if $\rank(L_A) = \dim(F^n) = n$. Thus we conclude that, an $n \times n$ matrix is invertible if and only if its rank is $n$.

\begin{theorem}
    Let $T:V \to W$ be a linear transformation between finite-dimensional vector spaces, and let $\B$ and $\mathscr{C}$ be ordered basis for $V$ and $W$, respectively. Then $\rank(T) = \rank([T]_\B^\mathscr{C})$.
\end{theorem}

\begin{proof}
    We know $\rank(T) = \rank(L_A)$ where $A = [T]_\B^\mathscr{C}$. Also from the above definition we have $\rank(A) = \rank(L_A)$. Thus, $\rank(T) = \rank(A) = \rank([T]_\B^\mathscr{C})$.
\end{proof}

\begin{theorem}
Let $A$ be an $m \times n$ matrix. If $P$ and $Q$ are invertible $m \times m$ and $n \times n$ matrices, respectively, then
\begin{enumerate}
    \item $\rank(AQ) = \rank(A)$
    \item $\rank(PA) = \rank(A)$
    \item $\rank(PAQ) = \rank(A)$
\end{enumerate}
\end{theorem}

\begin{proof}
Here,
\[
\rank(A) = \rank(L_A) \;; \quad L_A : F^n \to F^m
\]
Also since $P$ and $Q$ are invertible, then, we can say that $L_P$ and $L_Q$ are invertible, respectively and thus, $L_P$ and $L_Q$ are onto.

\begin{enumerate}
    \item $\rank(AQ) = \dim(R(L_{AQ}))$.\\
    Now,
    \[
    R(L_{AQ}) = R(L_A \cdot L_Q) = L_A L_Q(F^n) = L_A(L_Q(F^n)) = L_A(F^n) = R(L_A)
    \]
    Thus, $\rank(AQ) = \dim(R(L_{AQ})) = \dim(R(L_A)) = \rank(A)$.

    \item $\rank(PA) = \dim(R(L_{PA})) = \dim(L_P \cdot L_A(F^n)) = \dim(L_P(L_A(F^n))) = \dim(L_P(R(L_A)))$.\\
    Since $R(L_A)$ is a subspace of $F^m$, and we know that if $T$ is a transformation from $V$ to $W$ is an isomorphism and $V_0$ is a subspace of $V$ then $\dim(V_0) = \dim(T(V_0))$. Therefore $\dim(L_P(R(L_A))) = \dim(R(L_A))$. Thus,
    \[
    \rank(PA) = \dim(L_P(R(L_A))) = \dim(R(L_A)) = \rank(A).
    \]

    \item Using 1 and 2 simultaneously, we get,
    \[
    \rank(PAQ) = \rank(PA) = \rank(A).
    \]
\end{enumerate}
\end{proof}

\begin{corollary}
Elementary row and column operations on a matrix are rank-preserving.
\end{corollary}

\begin{proof}
If $B$ is obtained from a matrix $A$ by an elementary row operation, then there exists an elementary matrix $E$ such that $B = EA$. By Theorem 2.1, $E$ is invertible and by Theorem 3.4 (2),
\[
\rank(B) = \rank(EA) = \rank(A).
\]
Again, if $B$ is obtained from a matrix $A$ by an elementary column operation, then there exists an elementary matrix $E$ such that $B = AE$. By Theorem 3.2, $E$ is invertible and by Theorem 3.2 (1),
\[
\rank(B) = \rank(AE) = \rank(A).
\]
Thus we conclude that row and column operations are rank-preserving.
\end{proof}

\begin{theorem}
The rank of any matrix equals the maximum number of its linearly independent columns; that is, rank of a matrix is the dimension of the subspace generated by its columns.
\end{theorem}

\begin{proof}
For any $A \in M_{m \times n}(F)$, $\rank(A) = \rank(L_A) = \dim(R(L_A))$. Let $\B = \{e_1, e_2, \dots, e_n\}$ be the standard ordered basis for $F^n$. Then $\B$ spans $F^n$ and we know that for any transformation $T$ from finite-dimensional vector spaces $V$ with basis $\beta$ to $W$, $R(T) = \Span(T(\beta))$, therefore
\[
R(L_A) = \Span(L_A(\B)) = \Span(\{L_A(e_1), L_A(e_2), \dots, L_A(e_n)\}).
\]
We know, $L_A(e_j) = Ae_j$ and we know that, $Ae_j = a_j$ where $a_j$ is the $j$-th column of $A$. Hence,
\[
R(L_A) = \Span(\{a_1, a_2, \dots, a_n\}).
\]
Thus, \[\rank(A) = \dim(R(L_A)) = \dim(\Span(\{a_1, a_2, \dots, a_n\}))\].
\end{proof}

\begin{lemma}
For any matrix $A \in M_{m \times n}(F)$, $\rank(A) = 0$ iff $A = 0$.
\end{lemma}

\begin{proof}
Let $\rank(A) = 0$ then $\rank(A) = \rank(L_A) = 0 \;; \quad L_A : F^n \to F^m$.

From Theorem 3.3,
\[
\rank(A) = \dim(\Span(\{a_1, a_2, \dots, a_n\}))
\]
where $a_j$ ($1 \le j \le n$) is the $j$-th column of $A$. Thus,
\begin{align*}
& \dim(\Span(\{a_1, a_2, \dots, a_n\})) = 0 \\
& \Rightarrow \Span(\{a_1, a_2, \dots, a_n\}) = \{0\}.
\end{align*}
Thus, $a_1 = a_2 = \dots = a_n = 0$. Hence $A = 0$.

Conversely, let $A = 0$, then,
\[
A = \left. \underbrace{\begin{bmatrix}
0 & 0 & \dots & 0 \\
0 & 0 & \dots & 0 \\
\vdots & \vdots & \ddots & \vdots \\
0 & 0 & \dots & 0
\end{bmatrix}}_{n \text{ columns}} \right\} m \text{ rows}
\]

Now, from definition, $\rank(A) = \rank(L_A) \;; \quad L_A : F^n \to F^m$. Now, $\rank(L_A) = \dim(R(L_A))$. Since $L_A$ is a linear transformation then, $R(L_A) = L_A(F^n)$. By definition of $L_A$:
\[
L_A(F^n) = Ax \quad \forall \ x \in F^n.
\]
Since $A$ is a zero matrix, every vector in $F^n$ is reduced to $0$. Thus, $R(L_A) = L_A(F^n) = \{0\}$.

Therefore,
\[
\rank(A) = \rank(L_A) = \dim(R(L_A)) = 0.
\]
\end{proof}

\begin{lemma}
Let $B' \in M_{m \times n}(F)$ be a submatrix of $B \in M_{(m+1) \times (n+1)}(F)$ such that
\[
B = \left[
\begin{array}{c|ccc}
1 & 0 & \dots & 0 \\ \hline
0 &   &       &   \\
\vdots & & B'    &   \\
0 &   &       &   
\end{array}
\right]
\]
Then if $\rank(B) = \alpha$, then $\rank(B') = \alpha - 1$.
\end{lemma}

\begin{proof}
We know,
\[
\rank(B) = \alpha = \dim(R(L_B)) \;; \quad L_B : F^{n+1} \to F^{m+1}
\]
and
\[
\rank(B') = \dim(R(L_{B'})) \;; \quad L_{B'} : F^n \to F^m
\]
Now, Let,
\[
B' = \begin{bmatrix}
b'_{11} & b'_{12} & \dots & b'_{1n} \\
b'_{21} & b'_{22} & \dots & b'_{2n} \\
\vdots & \vdots & \ddots & \vdots \\
b'_{m1} & b'_{m2} & \dots & b'_{mn}
\end{bmatrix}_{m \times n}
\text{ and }
B = \begin{bmatrix}
1 & 0 & 0 & \dots & 0 \\
0 & b'_{11} & b'_{12} & \dots & b'_{1n} \\
0 & b'_{21} & b'_{22} & \dots & b'_{2n} \\
\vdots & \vdots & \vdots & \ddots & \vdots \\
0 & b'_{m1} & b'_{m2} & \dots & b'_{mn}
\end{bmatrix}_{(m+1) \times (n+1)}
\]
Let, $b_1, b_2, \dots, b_{n+1}$ denote the $(n+1)$ columns of $B$ and $b'_1, b'_2, \dots, b'_n$ denote the $n$ columns of $B'$. From Theorem 3.3, we have,
\[
\rank(B) = \dim(\Span(\{b_1, b_2, \dots, b_{n+1}\}))
\]
Since, $b_1$ is linearly independent from other columns in $B$. Therefore,
\[
\rank(B) = 1 + \dim(\Span(\{b_2, b_3, \dots, b_{n+1}\})).
\]
Here, in $B$, if $b_j$ denote the $j$-th column of $B$ then,
\[
b_j = \begin{pmatrix} 0 \\ b'_{j-1} \end{pmatrix}
\]
where $b'_j$ is the $j$-th column of $B'$ for $j \in \{2, 3, \dots, n+1\}$. This,
\begin{equation}
\rank(B) = 1 + \dim\left(\Span\left(\left\{ \begin{pmatrix} 0 \\ b'_1 \end{pmatrix}, \begin{pmatrix} 0 \\ b'_2 \end{pmatrix}, \dots, \begin{pmatrix} 0 \\ b'_n \end{pmatrix} \right\}\right)\right). \tag{1}
\end{equation}

Let $W = \left\{ \begin{pmatrix} 0 \\ x \end{pmatrix} \in F^{m+1} : x \in F^m \right\}$. Let us define
\[
\phi : W \to F^m \text{ by } \phi \begin{pmatrix} 0 \\ x \end{pmatrix} = x
\]
Again, let, $\begin{pmatrix} 0 \\ x \end{pmatrix}, \begin{pmatrix} 0 \\ y \end{pmatrix} \in W$ and $c \in F$ then,
\[
\phi \left[ c \begin{pmatrix} 0 \\ x \end{pmatrix} + \begin{pmatrix} 0 \\ y \end{pmatrix} \right] = \phi \left[ \begin{pmatrix} 0 \\ cx + y \end{pmatrix} \right] = cx + y = c \phi \begin{pmatrix} 0 \\ x \end{pmatrix} + \phi \begin{pmatrix} 0 \\ y \end{pmatrix}.
\]
Thus, $\phi$ is linear. Now, let $x \in N(\phi)$, then
\begin{align*}
& \phi \begin{pmatrix} 0 \\ x \end{pmatrix} = 0_{F^m} \\
& \Rightarrow x = 0_{F^m} \\
& \therefore N(\phi) = \begin{pmatrix} 0 \\ 0_{F^m} \end{pmatrix}
\end{align*}
Thus $\phi$ is one-to-one. Again since, $R(\phi)$ is a subspace of $F^m$ then $R(\phi) \subseteq F^m$. Also for any $y \in F^m$,
\begin{align*}
& \phi \begin{pmatrix} 0 \\ y \end{pmatrix} = y \in R(\phi) \\
& \therefore F^m \subseteq R(\phi).
\end{align*}
Thus $R(\phi) = F^m$ and therefore $\phi$ is onto and hence $\phi$ is an isomorphism. Thus $W$ is isomorphic to $F^m$.
\[
\therefore \dim\left(\Span\left(\left\{ \begin{pmatrix} 0 \\ b'_1 \end{pmatrix}, \begin{pmatrix} 0 \\ b'_2 \end{pmatrix}, \dots, \begin{pmatrix} 0 \\ b'_n \end{pmatrix} \right\}\right)\right) = \dim(\Span(\{b'_1, b'_2, \dots, b'_n\}))
\]
\begin{align*}
    \therefore \ (1) \implies \rank(B) &= 1 + \dim(\Span(\{b'_1, b'_2, \dots, b'_n\})) \\
    &= 1 + \rank(B') \quad \text{[using Theorem 3.3]} \\
    \implies \rank(B') &= \alpha - 1.
\end{align*}
\end{proof}

\begin{lemma}
Let $B'$ and $D'$ be $m \times n$ matrices and $B$ and $D$ be $(m+1) \times (n+1)$ matrices defined by
\[
B = \left[
\begin{array}{c|ccc}
1 & 0 & \dots & 0 \\ \hline
0 &   &       &   \\
\vdots & & B'    &   \\
0 &   &       &   
\end{array}
\right]
\quad \text{and} \quad
D = \left[
\begin{array}{c|ccc}
1 & 0 & \dots & 0 \\ \hline
0 &   &       &   \\
\vdots & & D'    &   \\
0 &   &       &   
\end{array}
\right].
\]
Then if $B'$ can be transformed into $D'$ by an elementary row/column operation, then $B$ can be transformed into $D$ by an elementary row/column operation.
\end{lemma}

\begin{proof}
\textbf{Proving for Elementary Row Operations:}\\
We have, $D'$ is obtained by an elementary row operation from $B'$ then $D' = E'B'$ where $E'$ is an $m \times m$ elementary matrix obtained from $I_m$.

We need to prove $D = EB$ where $E$ is an $(m+1) \times (m+1)$ elementary matrix obtained from $I_{m+1}$ and $E$ is defined by
\[
E = \left[
\begin{array}{c|ccc}
1 & 0 & \dots & 0 \\ \hline
0 &   &       &   \\
\vdots & & E'    &   \\
0 &   &       &   
\end{array}
\right].
\]
We can write the matrices as,
\begin{align*}
    D &= \begin{bmatrix} 1 & 0_{1 \times n} \\ 0_{m \times 1} & D' \end{bmatrix}, \quad 
B = \begin{bmatrix} 1 & 0_{1 \times n} \\ 0_{m \times 1} & B' \end{bmatrix} \quad \text{and} \quad 
E = \begin{bmatrix} 1 & 0_{1 \times n} \\ 0_{m \times 1} & E' \end{bmatrix} \\
&= \begin{bmatrix} 1 & 0_{1 \times n} \\ 0_{m \times 1} & E'B' \end{bmatrix}
\end{align*}
Now,
\begin{align*}
EB &= \begin{bmatrix} 1 & 0_{1 \times n} \\ 0_{m \times 1} & E' \end{bmatrix} \begin{bmatrix} 1 & 0_{1 \times n} \\ 0_{m \times 1} & B' \end{bmatrix} \\
&= \begin{bmatrix} 1 \times 1 + 0_{1 \times n} \times 0_{m \times 1} & 1 \times 0_{1 \times n} + 0_{1 \times n} \times B' \\ 0_{m \times 1} \times 1 + E' \times 0_{m \times 1} & 0_{m \times 1} \times 0_{1 \times n} + E'B' \end{bmatrix} \\
&= \begin{bmatrix} 1 & 0_{1 \times n} \\ 0_{m \times 1} & E'B' \end{bmatrix} \\
&= D.
\end{align*}

\noindent\textbf{Proving for Elementary Column Operations:}\\
We have $D'$ is obtained by an elementary column operation from $B'$ then $D' = B'E'$ where $E'$ is an $n \times n$ elementary matrix obtained from $I_n$. We need to prove $D = BE$ where $E$ is an $(n+1) \times (n+1)$ elementary matrix obtained from $I_{n+1}$ and $E$ is defined by
\[
E = \left[
\begin{array}{c|ccc}
1 & 0 & \dots & 0 \\ \hline
0 &   &       &   \\
\vdots & & E'    &   \\
0 &   &       &   
\end{array}
\right] = \begin{bmatrix} 1 & 0_{1 \times n} \\ 0_{m \times 1} & E' \end{bmatrix}.
\]
Now,
\begin{align*}
BE &= \begin{bmatrix} 1 & 0_{1 \times n} \\ 0_{m \times 1} & B' \end{bmatrix} \begin{bmatrix} 1 & 0_{1 \times n} \\ 0_{m \times 1} & E' \end{bmatrix} \\
&= \begin{bmatrix} 1 \times 1 + 0_{1 \times n} \times 0_{m \times 1} & 1 \times 0_{1 \times n} + 0_{1 \times n} \times E' \\ 0_{m \times 1} \times 1 + B' \times 0_{m \times 1} & 0_{m \times 1} \times 0_{1 \times n} + B'E' \end{bmatrix} \\
&= \begin{bmatrix} 1 & 0_{1 \times n} \\ 0_{m \times 1} & B'E' \end{bmatrix} \\
&= \begin{bmatrix} 1 & 0_{1 \times n} \\ 0_{m \times 1} & D' \end{bmatrix} \\
&= D.
\end{align*}
\end{proof}

\begin{theorem}
Let $A$ be an $m \times n$ matrix of rank $r$. Then $r \le m$, $r \le n$ and, by means of a finite number of elementary row and column operations, $A$ can be transformed into the matrix
\[
D = \begin{bmatrix} I_r & O_1 \\ O_2 & O_3 \end{bmatrix}
\]
where $O_1, O_2$ and $O_3$ are zero matrices. Thus $D_{ii} = 1$ for $i \le r$ and $D_{ij} = 0$ otherwise.
\end{theorem}

\begin{proof}
If $A = 0$ then by Lemma 3.1, we have $\rank(A) = 0$. Thus, $D$ becomes,
\[
D = \begin{bmatrix} I_0 & O_1 \\ O_2 & O_3 \end{bmatrix}
\]
which is also a zero matrix. Hence $D = A$.

Now suppose $A \ne 0$ and $\rank(A) = r$, then $r > 0$. We use mathematical induction to establish the theorem.

Suppose $m = 1$, we then have $A$ as
\[
A = \begin{bmatrix} A_{11} & A_{12} & \dots & A_{1n} \end{bmatrix}.
\]
Here we perform at most one type 1 column operation and at most one type 2 column operation and we transform $A$ to
\[
D = \begin{bmatrix} 1 & A_{12} & \dots & A_{1n} \end{bmatrix}.
\]
Now we perform at most $n-1$ type 3 column operations to get,
\[
D = \begin{bmatrix} 1 & 0 & \dots & 0 \end{bmatrix}.
\]
By corollary 3.2.1 and Theorem 3.3, we conclude
\[
\rank(A) = \rank(D) = 1.
\]

Now suppose $n = 1$ then similarly performing at most one type 1 and at most one type 2 row operations and at most $m-1$ type 3 row operations, $A$ is converted to
\[
D = \begin{bmatrix} 1 & 0 & \dots & 0 \end{bmatrix}^t.
\]
Again by Theorem 3.3 and corollary 3.2.1,
\[
\rank(A) = \rank(D) = 1.
\]
Thus the theorem holds true for $m = 1$ and $n = 1$.

Now let us take a mathematical hypothesis that the theorem holds true for $m-1$ rows and $n-1$ columns for some $m, n > 0$. Since $A \ne 0$ then $A_{ij} \ne 0$ for some $i, j$. Now by at most one type 1 row and column operations we can bring non-zero entry to $1,1$ position and then by at most one type 2 operation we can get $1$ in $1,1$ position. Now by at most $m-1$ type 3 row operations and at most $n-1$ type 3 column operations we can get $0$ in first row and first column except $1$ in $1,1$ position. Thus $A$ can be transformed into
\[
B = \left[
\begin{array}{c|ccc}
1 & 0 & \dots & 0 \\ \hline
0 &   &       &   \\
\vdots & & B'    &   \\
0 &   &       &   
\end{array}
\right]
\]
where $B'$ is an $(m-1) \times (n-1)$ matrix. By corollary 3.2.1,
\[
\rank(A) = \rank(B) = r.
\]
By Lemma 3.2, $\rank(B') = r - 1$. Thus by induction hypothesis,
\[
r - 1 \le m - 1 \quad \text{and} \quad r - 1 \le n - 1.
\]
Hence, $r \le m$ and $r \le n$.

Also by our hypothesis, $B'$ can be transformed by elementary operations into
\[
D' = \begin{bmatrix} I_{r-1} & O_4 \\ O_5 & O_6 \end{bmatrix}
\]
where $O_4, O_5$ and $O_6$ are zero matrices. Thus $D'$ consists of all zero entries except for $r-1$ diagonal entries, which are ones. Let,
\[
D = \left[
\begin{array}{c|ccc}
1 & 0 & \dots & 0 \\ \hline
0 &   &       &   \\
\vdots & & D'    &   \\
0 &   &       &   
\end{array}
\right]
\]
Since $B'$ can be transformed into $D'$ then by Lemma 3.3, $B$ can be transformed into $D$ by finite elementary operations. Since $A$ can be transformed into $B$ and $B$ can be transformed into $D$ by finite elementary operation then $A$ can be transformed into $D$ by finite elementary operations.

Now as $D'$ contains $r-1$ diagonal entries as ones then $D$ contains $r$ diagonal entries as ones. Thus we have an $m \times n$ matrix $A$ with rank $r$ then $r \le m$ and $r \le n$ and $A$ can be transformed using finite elementary operations into
\[
D = \begin{bmatrix} I_r & O_7 \\ O_8 & O_9 \end{bmatrix}
\]
where $O_7, O_8$ and $O_9$ are zero matrices. Thus the theorem is established.
\end{proof}

\begin{corollary}
Let $A$ be an $m \times n$ matrix of rank $r$. Then there exist invertible matrices $B$ and $C$ of order $m \times m$ and $n \times n$, respectively, such that $D = BAC$, where
\[
D = \begin{bmatrix} I_r & O_1 \\ O_2 & O_3 \end{bmatrix}
\]
is the $m \times n$ matrix in which $O_1, O_2$ and $O_3$ are zero matrices.
\end{corollary}

\begin{proof}
By Theorem 3.4, $A$ can be transformed by means of a finite number of elementary row and column operations into the matrix $D$. We can say that each time we perform an elementary operation we multiply the matrix $A$ by elementary matrices of order $m \times m$ $E_1, E_2, \dots, E_p$ and elementary matrices of order $n \times n$ $G_1, G_2, \dots, G_q$ such that
\[
D = E_p E_{p-1} \dots E_2 E_1 A G_1 G_2 \dots G_q.
\]
By Theorem 2.1, each $E_i$ and $G_j$ are invertible. Let $B = E_p E_{p-1} \dots E_1$ and $C = G_1 G_2 \dots G_q$. Then $B$ and $C$ are invertible and $D = BAC$.
\end{proof}

\begin{corollary}
Let $A$ be an $m \times n$ matrix. Then
\begin{enumerate}
    \item $\rank(A^t) = \rank(A)$
    \item The rank of any matrix equals the maximum number of its linearly independent rows; that is, the rank of a matrix is the dimension of the subspace generated by its rows.
    \item The rows and columns of any matrix generate subspaces of the same dimension, numerically equal to the rank of the matrix.
\end{enumerate}
\end{corollary}

\begin{proof}
\begin{enumerate}
    \item By Corollary 3.4.1, there exist invertible matrices $B$ and $C$ such that $D = BAC$ where $D$ satisfies the stated conditions of the corollary. Taking transposes, we get
    \[
    D^t = (BAC)^t = C^t A^t B^t
    \]
    Since $B$ and $C$ are invertible, then $B^t$ and $C^t$ are also invertible. Thus by Theorem 3.2 (3),
    \[
    \rank(A^t) = \rank(C^t A^t B^t) = \rank(D^t).
    \]
    Suppose $\rank(A) = r$. Then $D^t$ is an $n \times m$ matrix from the matrix $D$ in Corollary 3.4.1.
    \[
    D^t = \begin{bmatrix} I_r & O_1 \\ O_2 & O_3 \end{bmatrix}^t = \begin{bmatrix} I_r & O_2 \\ O_1 & O_3 \end{bmatrix}
    \]
    Thus, $\rank(D^t) = r$. Thus,
    \[
    \rank(A^t) = \rank(D^t) = r = \rank(A).
    \]

    \item Let $A$ be an $m \times n$ matrix with linearly independent rows. Then applying transpose on $A$ makes it an $n \times m$ matrix with linearly independent columns. Thus by Theorem 3.3,
    \[
    \rank(A^t) = \dim(\Span(\{a_1, a_2, \dots, a_n\}))
    \]
    where $a_1, a_2, \dots, a_n$ are linearly independent columns of $A^t$ which are again linearly independent rows of $A$. Also from part (1),
    \[
    \rank(A) = \dim(\Span(\{a_1, a_2, \dots, a_n\})).
    \]
    Thus, rank of any matrix equals the maximum number of its linearly independent rows.

    \item This follows directly from part-1, part-2 of this theorem and from Theorem 3.3.
\end{enumerate}
\end{proof}

\begin{corollary}
Every invertible matrix is a product of elementary matrices.
\end{corollary}

\begin{proof}
If $A$ is an invertible $n \times n$ matrix then $\rank(A) = n$. Hence the matrix $D$ in Corollary 3.4.1 equals $I_n$, and there exist invertible matrices $B$ and $C$ such that $I_n = BAC$.

In Corollary 3.4.1 proof, we defined $B$ and $C$ as $B = E_p E_{p-1} \dots E_1$ and $C = G_1 G_2 \dots G_q$ where $E_i$'s and $G_j$'s are elementary matrices.

Now,
\begin{align*}
I_n &= BAC \\
\implies B^{-1} I_n C^{-1} &= B^{-1} B A C C^{-1} \\
\implies B^{-1} C^{-1} &= I_n A I_n \\
&= A.
\end{align*}
Thus,
\begin{align*}
A = B^{-1} C^{-1} &= (E_p E_{p-1} \dots E_1)^{-1} (G_1 G_2 \dots G_q)^{-1} \\
&= E_1^{-1} E_2^{-1} \dots E_p^{-1} G_q^{-1} G_{q-1}^{-1} \dots G_1^{-1}.
\end{align*}
From Theorem 2.1, we know inverse of an elementary matrix is also an elementary matrix. Thus $A$ is a product of elementary matrices.
\end{proof}

\begin{theorem}
Let $T: V \to W$ and $U: W \to Z$ be linear transformations on finite-dimensional vector spaces $V, W$ and $Z$, and let $A$ and $B$ be matrices such that the product $A B$ is defined. Then,
\begin{enumerate}
    \item $\rank(UT) \le \rank(U)$
    \item $\rank(UT) \le \rank(T)$
    \item $\rank(AB) \le \rank(A)$
    \item $\rank(AB) \le \rank(B)$
\end{enumerate}
\end{theorem}

\begin{proof}
We prove in order 1, 3, 4 and 2.
\begin{enumerate}
    \item We know $R(T) \subseteq W$. Thus,
    \[
    R(UT) = UT(V) = U(T(V)) = U(R(T)) \subseteq U(W) = R(U).
    \]
    Hence,
    \[
    \rank(UT) = \dim(R(UT)) \le \dim(R(U)) = \rank(U).
    \]

    \item[\text{3.}] By 1,
    \[
    \rank(AB) = \rank(L_{AB}) = \rank(L_A L_B) \le \rank(L_A) = \rank(A).
    \]

    \item[\text{4.}] By 3 and corollary 3.4.2 (1),
    \[
    \rank(AB) = \rank((AB)^t) = \rank(B^t A^t) \le \rank(B^t) = \rank(B).
    \]

    \item[\text{2.}] Let $\alpha, \beta$ and $\gamma$ be ordered bases for $V, W$ and $Z$, respectively, and let $A' = [U]_\beta^\gamma$ and $B' = [T]_{\alpha}^\beta$. Then $A'B' = [UT]_\alpha^\gamma$. Hence by Theorem 3.1 and part-4,
    \[
    \rank(UT) = \rank(A'B') \le \rank(B') = \rank(T).
    \]
\end{enumerate}
\end{proof}

\section{The Inverse of a Matrix}

\begin{definition}
    Let $A$ and $B$ be $m \times n$ and $m \times p$ matrices respectively. By the augmented matrix we mean $(A|B)$, we mean the $m \times (n+p)$ matrix, that is, the matrix whose first $n$ columns are columns of $A$, and whose last $p$ columns are the columns of $B$.
\end{definition}

\begin{theorem}
If $A$ and $B$ are matrices having $n$ rows. Prove that
\[
M(A|B) = (MA|MB)
\]
for any $m \times n$ matrix $M$.
\end{theorem}

\begin{proof}
Let $A$ be an $n \times p$ matrix and $B$ be an $n \times m$ matrix. Then $(A|B)$ is an augmented matrix of order $n \times (p+m)$. We can decompose $(A|B)$ into a sum of two matrices by padding them with appropriate sized zero matrices. Then,
\[
(A|B) = [A_{n \times p} | 0_{n \times m}] + [0_{n \times p} | B_{n \times m}]
\]
Thus,
\begin{align*}
M(A|B) &= M\Big([A_{n \times p} | 0_{n \times m}] + [0_{n \times p} | B_{n \times m}]\Big) \\
&= M_{m \times n}[A_{n \times p} | 0_{n \times m}] + M_{m \times n}[0_{n \times p} | B_{n \times m}].
\end{align*}
Since, we know for matrices $U$ and $V$, $j$-th column of $UV$ is equal to the product of $U$ and $j$-th column of $V$. Thus, the matrix multiplication distributes the multiplication across columns in the augmented matrix.
\begin{align*}
M(A|B) &= [M_{m \times n} A_{n \times p} | M_{m \times n} 0_{n \times m}] + [M_{m \times n} 0_{n \times p} | M_{m \times n} B_{n \times m}] \\
&= [(MA)_{m \times p} | 0_{m \times m}] + [0_{m \times p} | (MB)_{m \times m}] = (MA|MB)_{m \times (p+m)}.
\end{align*}
\end{proof}

Let $A$ be an invertible $n \times n$ matrix, and consider an $n \times 2n$ augmented matrix $C = (A | I_n)$. By Theorem 4.1,
\[
A^{-1}C = A^{-1}(A | I_n) = (A^{-1}A | A^{-1}I_n) = (I_n | A^{-1}).
\]
By Corollary 3.4.3, $A^{-1}$ is the product of elementary matrices, say $A^{-1} = E_p E_{p-1} \dots E_1$. Thus our above equation becomes
\[
A^{-1}C = E_p E_{p-1} \dots E_1 (A | I_n) = (I_n | A^{-1}).
\]
If $A$ is an invertible $n \times n$ matrix then it is possible to transform the matrix $(A | I_n)$ into the matrix $(I_n | A^{-1})$ by means of a finite number of elementary row operations.

Also suppose that $A$ is an invertible matrix and that, for some $n \times n$ matrix $B$, the matrix $(A | I_n)$ can be transformed into the matrix $(I_n | B)$ by a finite number of elementary row operations. Let $E_1, E_2, \dots, E_p$ be the elementary matrices associated with these row operations. Then,
\[
E_p E_{p-1} \dots E_1 (A | I_n) = (I_n | B).
\]
Let, $M = E_p E_{p-1} \dots E_1$ then,
\[
M(A | I_n) = (MA | M) = (I_n | B). \qquad [\text{By Theorem 4.1}]
\]
Thus, $MA = I_n$ and $M = B$. It follows that $M = A^{-1}$ or $B = A^{-1}$. Thus if $A$ is an invertible $n \times n$ matrix, and the matrix $(A | I_n)$ is transformed into a matrix of the form $(I_n | B)$ by means of a finite number of elementary row operations, then $B = A^{-1}$.


\section{Conclusion}
Throughout this paper, we have systematically unpacked the algebraic skeleton of matrix rank and its invariance under elementary operations. By establishing that row and column operations do not alter the underlying column space's dimension, we bridge the gap between computational matrix manipulations and the invariant properties of linear operators. The derivation of the rank canonical form highlights that rank serves as the complete invariant for matrix equivalence, while the column-distribution property of matrix multiplication provides the formal justification for computing matrix inverses through row reduction on augmented systems.

\begin{thebibliography}{9}
    \bibitem{fis_linear_algebra}
    Friedberg, S. H., Insel, A. J., \& Spence, L. E. (2018). \textit{Linear Algebra} (5th ed.). Pearson.
\end{thebibliography}
\end{document}
